% Part: first-order-logic
% Chapter: syntax-and-semantics
% Section: subformulas

\documentclass[../../../include/open-logic-section]{subfiles}

\begin{document}

\olfileid{fol}{syn}{sbf}

\olsection{\printtoken{P}{subformula}}

\begin{explain}
কোনো প্রদত্ত !!{formula}-কে যে !!{formula}s “গড়ে তোলে”, সেগুলি নিয়ে
কথা বলা প্রায়ই উপযোগী। এগুলিকে তার \emph{!!{subformula}s} বলি।
প্রত্যেকটি !!{formula} নিজেই নিজের !!a{subformula}; $!A$ নিজে ছাড়া
$!A$-এর অন্য কোনো উপসূত্রকে \emph{প্রকৃত !!{subformula}} বলে।
\end{explain}

\begin{defn}[অব্যবহিত !!^{subformula}]
$!A$ যদি !!a{formula} হয়, তাহলে $!A$-এর \emph{অব্যবহিত
!!{subformula}s} নিচের নিয়মে আরোহীভাবে সংজ্ঞায়িত:
\begin{enumerate}
\item পরমাণু !!{formula}s-এর কোনো অব্যবহিত !!{subformula}s নেই।

\tagitem{prvNot}{\indcase{!A}{\lnot !B}{$\indfrm$-এর একমাত্র অব্যবহিত
    !!{subformula} হলো~$!B$।}}{}

\item \indcase{!A}{(!B \ast !C)}{$\indfrm$-এর অব্যবহিত
  !!{subformula}s হলো $!B$ ও $!C$ ($\ast$ হলো দ্বি-স্থানীয়
  সংযোজকগুলির যেকোনো একটি)।}

\tagitem{prvAll}{\indcase{!A}{\lforall[x][!B]}{$\indfrm$-এর একমাত্র
    অব্যবহিত !!{subformula} হলো~$!B$।}}{}

\tagitem{prvEx}{\indcase{!A}{\lexists[x][!B]}{$\indfrm$-এর একমাত্র
    অব্যবহিত !!{subformula} হলো~$!B$।}}{}
\end{enumerate}
\end{defn}

\begin{defn}[প্রকৃত !!^{subformula}]
$!A$ যদি !!a{formula} হয়, তাহলে $!A$-এর \emph{প্রকৃত
!!{subformula}s} নিচের নিয়মে পুনরাবৃত্তিমূলকভাবে সংজ্ঞায়িত:
\begin{enumerate}
\item পরমাণু !!{formula}s-এর কোনো প্রকৃত !!{subformula}s নেই।

\tagitem{prvNot}{\indcase{!A}{\lnot !B}{$\indfrm$-এর প্রকৃত
    !!{subformula}s হলো $!B$ এবং $!B$-এর সব প্রকৃত
    !!{subformula}s।}}{}

\item \indcase{!A}{(!B \ast !C)}{$\indfrm$-এর প্রকৃত
  !!{subformula}s হলো $!B$, $!C$, $!B$-এর সব প্রকৃত !!{subformula}s
  এবং $!C$-এর সেই উপসূত্রগুলি।}

\tagitem{prvAll}{\indcase{!A}{\lforall[x][!B]}{$\indfrm$-এর প্রকৃত
    !!{subformula}s হলো $!B$ এবং $!B$-এর সব প্রকৃত
    !!{subformula}s।}}{}

\tagitem{prvEx}{\indcase{!A}{\lexists[x][!B]}{$\indfrm$-এর প্রকৃত
    !!{subformula}s হলো $!B$ এবং $!B$-এর সব প্রকৃত
    !!{subformula}s।}}{}
\end{enumerate}
\end{defn}

\begin{defn}[!!^{subformula}]
$!A$-এর !!{subformula}s হলো $!A$ নিজে এবং তার সব প্রকৃত
!!{subformula}s।
\end{defn}

\begin{explain}
অব্যবহিত !!{subformula}s ও প্রকৃত !!{subformula}s-কে আমরা যেভাবে
সংজ্ঞায়িত করেছি, তার সূক্ষ্ম পার্থক্যটি লক্ষ করুন। প্রথম ক্ষেত্রে
$!A$-এর সম্ভাব্য প্রতিটি রূপের জন্য সরাসরি $!A$-এর অব্যবহিত
!!{subformula}s সংজ্ঞায়িত করেছি। এটি ক্ষেত্রভিত্তিক একটি স্পষ্ট
সংজ্ঞা, আর ক্ষেত্রগুলি !!{formula}s-এর সেটের আরোহী সংজ্ঞাকে অনুসরণ
করে। দ্বিতীয় ক্ষেত্রেও সব !!{formula}s-এর সেট যেভাবে সংজ্ঞায়িত,
আমরা তা অনুসরণ করেছি; কিন্তু প্রতিটি ক্ষেত্রে ছোট !!{formula}s $!B$,
$!C$-এর প্রকৃত !!{subformula}s নেওয়ার পাশাপাশি এই !!{formula}s-কেও
অন্তর্ভুক্ত করেছি। এতে সংজ্ঞাটি \emph{পুনরাবৃত্তিমূলক} হয়। সাধারণভাবে, আরোহীভাবে
সংজ্ঞায়িত কোনো সেটে (এখানে !!{formula}s) কোনো ফলনের সংজ্ঞাকে
পুনরাবৃত্তিমূলক বলা হয়, যদি ফলনের সংজ্ঞার ক্ষেত্রগুলি সেই ফলনকেই
ব্যবহার করে। তবে সংজ্ঞাটি সুসংজ্ঞায়িত হতে হলে নিশ্চিত করতে হবে যে,
আমরা শুধু যে যুক্তিকে সংজ্ঞায়িত করছি তার “আগে” আসা যুক্তিগুলির জন্যই
ফলনটির মান ব্যবহার করি। এখানে $(!B \ast !C)$-এর “প্রকৃত
!!{subformula}” সংজ্ঞায়িত করার সময় কেবল “আগের” !!{formula}s $!B$ ও
$!C$-এর প্রকৃত !!{subformula}s ব্যবহার করি।
\end{explain}

\begin{prop}
\ollabel{prop:subfrm-trans}
ধরি $!B$ হলো $!A$-এর একটি উপসূত্র এবং $!C$ হলো $!B$-এর একটি উপসূত্র।
তাহলে $!C$ হলো $!A$-এর একটি উপসূত্র। অন্যভাবে বললে, উপসূত্র-সম্বন্ধটি
সঞ্চারী।
\end{prop}

\begin{prob}
\olref[fol][syn][sbf]{prop:subfrm-trans} প্রমাণ করুন।
\end{prob}

\begin{prop}
\ollabel{prop:count-subfrms}
ধরি $!A$ এমন একটি সূত্র যাতে $n$টি সংযোজক ও পরিমাণসূচক আছে।
তাহলে $!A$-এর উপসূত্রের সংখ্যা সর্বাধিক $2n+1$।
\end{prop}

\begin{prob}
\olref[fol][syn][sbf]{prop:count-subfrms} প্রমাণ করুন।
\end{prob}

\end{document}
